\documentclass[aspectratio=169,11pt]{beamer}

% ============================================
% PACKAGES
% ============================================
\usepackage[utf8]{inputenc}
\usepackage[T1]{fontenc}
\usepackage[ngerman]{babel}
\usepackage{lmodern}
\usepackage{tcolorbox}
\tcbuselibrary{skins, hooks}
\usepackage{listings}
\usepackage{xcolor}
\usepackage{fontawesome5}
\usepackage{tikz}
\usetikzlibrary{shapes, arrows.meta, positioning}

% Modern Font
\IfFileExists{FiraSans.sty}{\usepackage[sfdefault]{FiraSans}}{}
\renewcommand{\familydefault}{\sfdefault}

% ============================================
% COLORS & DESIGN SYSTEM
% ============================================
\definecolor{BgColor}{HTML}{FAFAFA}
\definecolor{Primary}{HTML}{2B3A42}
\definecolor{Accent}{HTML}{E84A5F}
\definecolor{Secondary}{HTML}{3F5765}
\definecolor{CodeBg}{HTML}{F0F0F0}
\definecolor{Success}{HTML}{28A745}

\setbeamercolor{background canvas}{bg=BgColor}
\setbeamercolor{title}{fg=Primary}
\setbeamercolor{frametitle}{bg=Primary, fg=white}
\setbeamercolor{structure}{fg=Accent}
\setbeamercolor{normal text}{fg=Primary}
\setbeamercolor{itemize item}{fg=Accent}

\setbeamertemplate{navigation symbols}{}
\setbeamertemplate{footline}{%
  \leavevmode%
  \hbox{%
  \begin{beamercolorbox}[wd=.333333\paperwidth,ht=2.25ex,dp=1ex,center]{author in head/foot}%
    \usebeamerfont{author in head/foot}\tiny Falko, Malte, Max, Jana
  \end{beamercolorbox}%
  \begin{beamercolorbox}[wd=.333333\paperwidth,ht=2.25ex,dp=1ex,center]{title in head/foot}%
    \usebeamerfont{title in head/foot}\tiny Grundlagen der Praktischen Informatik
  \end{beamercolorbox}%
  \begin{beamercolorbox}[wd=.333333\paperwidth,ht=2.25ex,dp=1ex,right]{date in head/foot}%
    \usebeamerfont{date in head/foot}\insertframenumber{} / \inserttotalframenumber\hspace*{2ex} 
  \end{beamercolorbox}}%
  \vskip0pt%
}

% ============================================
% CUSTOM BOXES
% ============================================
\tcbset{
    modernbox/.style={
        enhanced,
        boxrule=0pt,
        frame hidden,
        colback=white,
        coltitle=white,
        colbacktitle=Secondary,
        fonttitle=\bfseries,
        drop lifted shadow,
        arc=4pt,
        toptitle=2mm,
        bottomtitle=2mm,
        left=2mm,
        right=2mm
    },
    successbox/.style={
        modernbox,
        colbacktitle=Success
    }
}

% ============================================
% CODE LISTINGS
% ============================================
\lstdefinelanguage{Haskell}{
  keywords={class, data, deriving, do, else, if, in, import, let, module, of, then, type, where, case, error},
  keywordstyle=\color{Accent}\bfseries,
  identifierstyle=\color{Primary},
  sensitive=true,
  comment=[l]{--},
  commentstyle=\color{gray}\ttfamily\itshape,
  stringstyle=\color{teal}\ttfamily,
  morestring=[b]"
}

\lstset{
    language=Haskell,
    backgroundcolor=\color{CodeBg},
    basicstyle=\ttfamily\scriptsize,
    breaklines=true,
    frame=single,
    rulecolor=\color{CodeBg},
    frameround=tttt,
    numbers=left,
    numberstyle=\tiny\color{gray},
    xleftmargin=15pt,
    tabsize=4,
    showstringspaces=false
}

% ============================================
% TITLE PAGE
% ============================================
\title{\textbf{Tutorium 5}}
\subtitle{Listen \& Listenkomprehension}
\institute{Grundlagen der Praktischen Informatik}
\author{}
\date{\today}

\begin{document}

\begin{frame}[plain]
    \titlepage
\end{frame}

\begin{frame}{Inhaltsverzeichnis}
    \tableofcontents
\end{frame}

% ============================================
% LISTEN
% ============================================
\section{Listen in Haskell}

\begin{frame}[fragile]{Was ist eine Liste?}
    Eine Liste ist ein algebraischer Datentyp, der rekursiv definiert ist:
    \[ \texttt{data [a] = [] | a : [a]} \]
    
    \begin{itemize}
        \item \texttt{[]} ist die \textbf{leere Liste} (Basisfall).
        \item \texttt{a : [a]} ist eine Liste mit einem \textbf{Kopf (Head)} und einem \textbf{Rest (Tail)}.
    \end{itemize}
    
    \begin{tcolorbox}[modernbox, title=Beispiel: Listenaufbau]
\begin{lstlisting}
-- Listenliteral als Syntax Sugar
[1, 2, 3, 4] == 1 : 2 : 3 : 4 : []

-- Pattern Matching auf Listen
head' :: [a] -> a
head' (x:xs) = x
\end{lstlisting}
    \end{tcolorbox}
    \small \textbf{Zusatz:} Strings sind einfach Listen von Characters: \texttt{\char`\"Hallo\char`\" == ['H', 'a', 'l', 'l', 'o']}.
\end{frame}

\begin{frame}[fragile]{Visualisierung der Listenstruktur}
    Eine Liste besteht aus aneinandergeketteten Elementen (Nodes).
    
    \begin{columns}[T]
        \begin{column}{0.48\textwidth}
            \textbf{Die Liste \texttt{[1, 2, 3]}}
            \vspace{0.3cm}
            \begin{center}
            \begin{tikzpicture}[node distance=0cm, every node/.style={draw=Primary, fill=BgColor, minimum height=0.6cm, font=\small}]
                % First node
                \node[minimum width=0.6cm, fill=CodeBg] (val1) {1};
                \node[minimum width=0.6cm, right=of val1] (ptr1) {$\bullet$};
                
                % Second node
                \node[minimum width=0.6cm, right=0.5cm of ptr1, fill=CodeBg] (val2) {2};
                \node[minimum width=0.6cm, right=of val2] (ptr2) {$\bullet$};
                
                % Third node
                \node[minimum width=0.6cm, right=0.5cm of ptr2, fill=CodeBg] (val3) {3};
                \node[minimum width=0.6cm, right=of val3] (ptr3) {$\bullet$};
                
                % End
                \node[draw=none, fill=none, right=0.5cm of ptr3] (end) {\texttt{[]}};
                
                \draw[-{Latex}, thick, Accent] (ptr1.center) -- (val2.west);
                \draw[-{Latex}, thick, Accent] (ptr2.center) -- (val3.west);
                \draw[-{Latex}, thick, Accent] (ptr3.center) -- (end.west);
            \end{tikzpicture}
            \end{center}
        \end{column}
        \begin{column}{0.48\textwidth}
            \textbf{Pattern Matching auf \texttt{(x:xs)}}
            \vspace{0.3cm}
            \begin{center}
            \begin{tikzpicture}[node distance=0cm, every node/.style={draw=Primary, fill=BgColor, minimum height=0.6cm, font=\small}]
                % First node
                \node[minimum width=0.6cm, fill=CodeBg] (val1) {1};
                \node[minimum width=0.6cm, right=of val1] (ptr1) {$\bullet$};
                
                % XS Box
                \node[draw=Secondary, fill=BgColor!90!Secondary, right=0.6cm of ptr1, minimum width=2.0cm, minimum height=0.8cm, dashed] (xs) {\texttt{[2, 3]}};
                
                \draw[-{Latex}, thick, Accent] (ptr1.center) -- (xs.west);
                
                % Labels
                \node[draw=none, fill=none, above=0.2cm of val1, font=\bfseries\color{Success}] {\texttt{x}};
                \node[draw=none, fill=none, above=0.2cm of xs, font=\bfseries\color{Success}] {\texttt{xs}};
            \end{tikzpicture}
            \end{center}
        \end{column}
    \end{columns}
\end{frame}

\begin{frame}[fragile]{Operationen auf Listen}
    Da Listen rekursiv sind, werden Operationen meist über \textbf{Pattern Matching} und Rekursion gelöst.
    
    \begin{columns}[T]
        \begin{column}{0.48\textwidth}
            \begin{tcolorbox}[modernbox, title=Element anfügen (Cons)]
\begin{lstlisting}
cons :: [a] -> a -> [a]
cons [] el     = [el]
cons (x:xs) el = x : cons xs el
\end{lstlisting}
            \end{tcolorbox}
        \end{column}
        \begin{column}{0.48\textwidth}
            \begin{tcolorbox}[modernbox, title=n-tes Element finden]
\begin{lstlisting}
nth :: Int -> [a] -> a
nth 0 (x:_)  = x
nth _ []     = error "Out of bounds"
nth n (_:xs) = nth (n-1) xs
\end{lstlisting}
            \end{tcolorbox}
        \end{column}
    \end{columns}
\end{frame}

% ============================================
% LISTENKOMPREHENSION
% ============================================
\section{Listenkomprehension}

\begin{frame}[fragile]{Listenkomprehension (List Comprehension)}
    Kompakte Schreibweise, um neue Listen aus alten zu erzeugen (ähnlich der Mengenlehre in der Mathematik).
    
    \textbf{Struktur:} \texttt{[ Ausdruck | Generator, Filter ]}
    
    \begin{tcolorbox}[successbox, title=Beispiele]
\begin{lstlisting}
-- Alle Quadratzahlen von 1 bis 10
squares = [x^2 | x <- [1..10]]

-- Nur gerade Quadratzahlen (mit Filter)
evenSquares = [x^2 | x <- [1..10], even x]

-- Vielfache von 3
multiplesOfThree = [x | x <- [1..30], x `mod` 3 == 0]
\end{lstlisting}
    \end{tcolorbox}
\end{frame}

% ============================================
% LIVE DEMO
% ============================================
\section{Praxis}

\begin{frame}[plain]
    \begin{center}
        \Huge \textbf{Live Demo} \\
        \vspace{0.5cm}
        \large \faLaptopCode~ \texttt{Listen.hs} \& \texttt{Listenkomprehension.hs} \\
        \vspace{0.3cm}
        \normalsize Gemeinsame Programmierung von Listenoperationen.
    \end{center}
    
    \vspace{0.5cm}
    \textbf{Aufgaben:}
    \begin{enumerate}
        \item Ein Element in einer Liste suchen (\texttt{findFirst}, \texttt{firstGreaterThan})
        \item Elemente an einem Index einfügen (\texttt{insertAt})
        \item Gefilterte Listen (z.\,B. \texttt{removeTwos})
    \end{enumerate}
\end{frame}

\end{document}
